\begin{table}[t]
\centering
\caption{Age-block sign diagnostic for $\log \Lambda-\varepsilon$ in the baseline benchmark}
\label{tab:baseline-net-term-sign}
\footnotesize
\renewcommand{\arraystretch}{1.08}
\begin{tabular}{lccccc}
\toprule
Age block & On-support share & Positive & Near zero & Negative & Mean $\log\Lambda-\varepsilon$ (plp) \\
\midrule
20-29 & 1.0000 & 0.2893 & 0.7020 & 0.0087 & 2.27 \\
30-39 & 1.0000 & 0.1532 & 0.8070 & 0.0398 & 1.14 \\
40-49 & 1.0000 & 0.1234 & 0.6161 & 0.2605 & 0.88 \\
50-59 & 1.0000 & 0.0932 & 0.2305 & 0.6763 & 0.49 \\
60-69 & 1.0000 & 0.0000 & 0.3259 & 0.6741 & -0.19 \\
70-79 & 1.0000 & 0.0672 & 0.2648 & 0.6680 & 0.23 \\
80-84 & 1.0000 & 0.5651 & 0.0000 & 0.4349 & 4.06 \\
\bottomrule
\end{tabular}
\begin{minipage}{0.90\textwidth}
\footnotesize Notes. ``On-support share'' is the share of younger mass in the age block that lies on one-step overlap support under the verified current-cell comparison map. Positive, near-zero, and negative columns are shares of on-support younger mass classified by the sign of $\log\Lambda-\varepsilon$. Values in the final column are reported in percent log points.
\end{minipage}
\end{table}
